\chapter{Contents}

\section{Contentable Spaces}

\begin{definition}[Contentable Space]
    A \textbf{contentable space} is a pair $(X, \mathcal{A})$ where $X$ is a set 
    and $\mathcal{A}$ is an algebra on $X$.
\end{definition}

\begin{intuition}
    A contentable space is the natural domain for defining finitely additive measures (contents). 
    The algebra structure provides closure under finite set operations, which is sufficient 
    for defining finite additivity. This is analogous to how a measurable space 
    (set + $\sigma$-algebra) is the natural domain for $\sigma$-additive measures.
\end{intuition}

\medskip
\begin{definition}[Contentable Subspace]
    Let $(X, \mathcal{A})$ be a contentable space and $Y \subseteq X$. 
    The \textbf{contentable subspace} induced on $Y$ is $(Y, \mathcal{A}_Y)$ where:
    \[
        \mathcal{A}_Y = \{ A \cap Y \mid A \in \mathcal{A} \}
    \]
\end{definition}

\begin{proposition}
    If $(X, \mathcal{A})$ is a contentable space and $Y \subseteq X$, then 
    $\mathcal{A}_Y$ is an algebra on $Y$.
\end{proposition}

\begin{proof}
    \textbf{Contains $\emptyset$:} $\emptyset = \emptyset \cap Y \in \mathcal{A}_Y$.

    \textbf{Contains $Y$:} $Y = X \cap Y \in \mathcal{A}_Y$.

    \textbf{Closure under complements:} Let $A \cap Y \in \mathcal{A}_Y$ with $A \in \mathcal{A}$. 
    Then $Y \setminus (A \cap Y) = (X \setminus A) \cap Y$. 
    Since $\mathcal{A}$ is an algebra, $X \setminus A \in \mathcal{A}$, so 
    $(X \setminus A) \cap Y \in \mathcal{A}_Y$.

    \textbf{Closure under unions:} Let $A_1 \cap Y, A_2 \cap Y \in \mathcal{A}_Y$. 
    Then $(A_1 \cap Y) \cup (A_2 \cap Y) = (A_1 \cup A_2) \cap Y \in \mathcal{A}_Y$.
\end{proof}

\bigskip
\section{Definition of Content}

\begin{definition}[Content / Premeasure]
    Let $(X, \mathcal{A})$ be a contentable space. A function $\mu: \mathcal{A} \to [0, \infty]$ 
    is a \textbf{content} (also called a \textbf{finitely additive measure} or \textbf{premeasure}) if:
    \begin{enumerate}
        \item $\mu(\emptyset) = 0$
        \item \textbf{Finite additivity:} If $A, B \in \mathcal{A}$ are disjoint, then 
            $\mu(A \cup B) = \mu(A) + \mu(B)$
    \end{enumerate}
\end{definition}

\begin{intuition}
    A content assigns ``sizes'' to sets in a way that respects finite disjoint unions. 
    This is the precursor to a measure, but lacks the crucial $\sigma$-additivity needed 
    for countable operations. Contents arise naturally:
    \begin{itemize}
        \item Jordan content on elementary sets measures ``nice'' bounded subsets of $\bb{R}^n$.
        \item Every measure restricted to an algebra is a content.
    \end{itemize}
    The properties of contents (monotonicity, subadditivity) follow from finite additivity alone.
\end{intuition}


\begin{proposition}[Properties of Contents]
Let $\mu$ be a content on an algebra $\mathcal{A}$. Then:
\begin{enumerate}
    \item \textbf{Monotonicity:} $A \subseteq B \implies \mu(A) \leq \mu(B)$
    \item \textbf{Subtractivity:} If $A \subseteq B$ and $\mu(A) < \infty$, then $\mu(B \setminus A) = \mu(B) - \mu(A)$
    \item \textbf{Finite subadditivity:} $\mu(A \cup B) \leq \mu(A) + \mu(B)$
\end{enumerate}
\end{proposition}

\begin{proof}
\textbf{Property 1 (Monotonicity):} Let $A \subseteq B$. We can write $B$ as a disjoint union:
\[ B = A \cup (B \setminus A) \]
By finite additivity: $\mu(B) = \mu(A) + \mu(B \setminus A)$. Since $\mu(B \setminus A) \geq 0$, we have $\mu(B) \geq \mu(A)$.

\medskip
\textbf{Property 2 (Subtractivity):} Let $A \subseteq B$ with $\mu(A) < \infty$. From the disjoint decomposition $B = A \cup (B \setminus A)$:
\[ \mu(B) = \mu(A) + \mu(B \setminus A) \]
Since $\mu(A) < \infty$, we can subtract: $\mu(B \setminus A) = \mu(B) - \mu(A)$.

\medskip
\textbf{Property 3 (Finite subadditivity):} Write $A \cup B$ as the disjoint union:
\[ A \cup B = A \cup (B \setminus A) \]
By finite additivity: $\mu(A \cup B) = \mu(A) + \mu(B \setminus A)$. By monotonicity (since $B \setminus A \subseteq B$): $\mu(B \setminus A) \leq \mu(B)$. Thus:
\[ \mu(A \cup B) \leq \mu(A) + \mu(B) \]
\end{proof}

\section{Jordan Content on $\bb{R}^n$}

\begin{goal}
Our goal is to define a notion of ``size'' or ``content'' for bounded subsets of $\bb{R}^n$ that:
\begin{itemize}
    \item Assigns rectangles their expected volume
    \item Is finitely additive on disjoint sets
    \item Can measure sets more general than rectangles
\end{itemize}
The Jordan content accomplishes this for ``nice'' sets---those whose boundary is negligible. This motivates the later development of Lebesgue measure, which can handle more pathological sets.
\end{goal}

\begin{definition}[Rectangle]
A \textbf{rectangle} (or \textbf{box}) in $\bb{R}^n$ is a set of the form:
\[ R = [a_1, b_1] \times [a_2, b_2] \times \cdots \times [a_n, b_n] \]
The \textbf{volume} of $R$ is $\text{vol}(R) = \prod_{i=1}^n (b_i - a_i)$.
\end{definition}

\begin{definition}[Elementary Set]
An \textbf{elementary set} is a finite union of rectangles. The collection of elementary sets is denoted $\mathcal{E}$.
\end{definition}

\begin{proposition}[Elementary Sets Form an Algebra]
$\mathcal{E}$ is an algebra on $\bb{R}^n$ (within any fixed bounding rectangle).
\end{proposition}

\begin{proof}
\textbf{Contains $\emptyset$:} The empty union is $\emptyset \in \mathcal{E}$.

\textbf{Closure under complements:} Let $E = \bigcup_{i=1}^k R_i$ be an elementary set. Then:
\[ E^c = \left(\bigcup_{i=1}^k R_i\right)^c = \bigcap_{i=1}^k R_i^c \]
Each $R_i^c$ (within a bounding box) is a finite union of rectangles (obtained by removing a rectangle from a box). The intersection of finitely many such sets is again a finite union of rectangles.

\textbf{Closure under finite unions:} By definition, finite unions of elementary sets are finite unions of rectangles, hence elementary.
\end{proof}

\begin{proposition}
Every elementary set can be written as a finite disjoint union of rectangles.
\end{proposition}

\begin{proof}
Proceed by induction on the number of rectangles in the union. Base case is trivial.

For the inductive step, let $E = R \cup E'$ where $E'$ is a disjoint union of rectangles $R_1, \ldots, R_k$.

Decompose $R$ into disjoint pieces: $R = (R \setminus \bigcup R_i) \cup \bigcup (R \cap R_i)$.

Each $R \cap R_i$ is a rectangle (intersection of rectangles). The set $R \setminus \bigcup R_i$ can be written as a finite disjoint union of rectangles by slicing along the faces of each $R_i$.
\end{proof}

\begin{definition}[Jordan Content]
The \textbf{Jordan content} $c$ on $\mathcal{E}$ is defined by:
\[ c(E) = \sum_{i=1}^k \text{vol}(R_i) \]
where $E = \bigcup_{i=1}^k R_i$ is a disjoint decomposition into rectangles.
\end{definition}

\begin{proposition}
Jordan content is well-defined and is a content on $\mathcal{E}$.
\end{proposition}

\begin{proof}
\textbf{Well-defined:} If $E = \bigcup R_i = \bigcup S_j$ are two disjoint decompositions, then $R_i = \bigcup_j (R_i \cap S_j)$. Since volume is additive on disjoint rectangles:
\[ \sum \text{vol}(R_i) = \sum_i \sum_j \text{vol}(R_i \cap S_j) = \sum_j \sum_i \text{vol}(R_i \cap S_j) = \sum \text{vol}(S_j) \]

\textbf{Content:} $c(\emptyset) = 0$ (empty sum). For disjoint $E, F \in \mathcal{E}$, write $E = \bigcup R_i$, $F = \bigcup S_j$. Then $E \cup F = (\bigcup R_i) \cup (\bigcup S_j)$, so $c(E \cup F) = c(E) + c(F)$.
\end{proof}

\section{Inner and Outer Content}

\begin{definition}[Outer Content (Axiomatic)]
    An \textbf{outer content} on $\Powerset{X}$ is a function $\mu^*: \Powerset{X} \to [0, \infty]$ satisfying:
    \begin{enumerate}
        \item[(O1)] $\mu^*(\emptyset) = 0$
        \item[(O2)] \textbf{Monotonicity:} $A \subseteq B \implies \mu^*(A) \leq \mu^*(B)$
        \item[(O3)] \textbf{Finite subadditivity:} $\mu^*(A \cup B) \leq \mu^*(A) + \mu^*(B)$
    \end{enumerate}
\end{definition}

\begin{proposition}[Subadditivity on Disjoint Sets]
For an outer content $\mu^*$, the following are equivalent:
\begin{enumerate}
    \item $\mu^*(A \cup B) \leq \mu^*(A) + \mu^*(B)$ for all $A, B$
    \item $\mu^*(A \cup B) \leq \mu^*(A) + \mu^*(B)$ for all disjoint $A, B$
\end{enumerate}
\end{proposition}

\begin{proof}
$(1) \Rightarrow (2)$: Immediate.

$(2) \Rightarrow (1)$: Let $A, B$ be arbitrary. Then $A \cup B = A \cup (B \setminus A)$ where $A$ and $B \setminus A$ are disjoint. By hypothesis:
\[ \mu^*(A \cup B) = \mu^*(A \cup (B \setminus A)) \leq \mu^*(A) + \mu^*(B \setminus A) \]
By monotonicity, $\mu^*(B \setminus A) \leq \mu^*(B)$. Thus $\mu^*(A \cup B) \leq \mu^*(A) + \mu^*(B)$.
\end{proof}

\bigskip
\begin{definition}[Inner Content (Axiomatic)]
    An \textbf{inner content} on $\Powerset{X}$ is a function $\mu_*: \Powerset{X} \to [0, \infty]$ satisfying:
    \begin{enumerate}
        \item[(I1)] $\mu_*(\emptyset) = 0$
        \item[(I2)] \textbf{Monotonicity:} $A \subseteq B \implies \mu_*(A) \leq \mu_*(B)$
        \item[(I3)] \textbf{Finite superadditivity on disjoint sets:} If $A \cap B = \emptyset$, then 
            $\mu_*(A \cup B) \geq \mu_*(A) + \mu_*(B)$
    \end{enumerate}
\end{definition}

\begin{remark}
The asymmetry between outer and inner content reflects their definitions via infima and suprema:
\begin{itemize}
    \item Outer content uses $\inf$ of covers, leading to subadditivity: covering $A \cup B$ separately may use more total volume.
    \item Inner content uses $\sup$ of inscribed sets, leading to superadditivity: disjoint inscribed sets can be combined.
\end{itemize}
\end{remark}

\subsection{Jordan Inner and Outer Content}

\begin{definition}[Jordan Outer Content]
    For any bounded $A \subseteq \bb{R}^n$, the \textbf{Jordan outer content} is:
    \[
        \JordanOuter{A} = \inf \{ \Jordan{E} \mid E \in \mathcal{E}, A \subseteq E \}
    \]
\end{definition}

\medskip
\begin{definition}[Jordan Inner Content]
    For any bounded $A \subseteq \bb{R}^n$, the \textbf{Jordan inner content} is:
    \[
        \JordanInner{A} = \sup \{ \Jordan{E} \mid E \in \mathcal{E}, E \subseteq A \}
    \]
\end{definition}

\bigskip
\begin{proposition}[Jordan Contents Satisfy the Axioms]
    The Jordan outer content $\JordanOuter{\cdot}$ satisfies (O1)--(O3), and the Jordan inner content $\JordanInner{\cdot}$ satisfies (I1)--(I3).
\end{proposition}

\begin{proof}
    \textbf{Outer content axioms:}

    \textbf{(O1)} $\JordanOuter{\emptyset} = \inf\{\Jordan{E} \mid \emptyset \subseteq E\} = \Jordan{\emptyset} = 0$.

    \textbf{(O2)} If $A \subseteq B$, then any cover of $B$ is a cover of $A$. 
    Thus $\JordanOuter{A} \leq \JordanOuter{B}$.

    \textbf{(O3)} If $A \subseteq E_1$ and $B \subseteq E_2$, then $A \cup B \subseteq E_1 \cup E_2$ and 
    $\Jordan{E_1 \cup E_2} \leq \Jordan{E_1} + \Jordan{E_2}$ (finite subadditivity of Jordan content). 
    Taking infima gives $\JordanOuter{A \cup B} \leq \JordanOuter{A} + \JordanOuter{B}$.

    \bigskip
    \textbf{Inner content axioms:}

    \textbf{(I1)} $\JordanInner{\emptyset} = \sup\{\Jordan{E} \mid E \subseteq \emptyset\} = \Jordan{\emptyset} = 0$.

    \textbf{(I2)} If $A \subseteq B$, then any set contained in $A$ is also contained in $B$. 
    Thus $\JordanInner{A} \leq \JordanInner{B}$.

    \textbf{(I3)} Let $A \cap B = \emptyset$. If $E_1 \subseteq A$ and $E_2 \subseteq B$, then 
    $E_1 \cap E_2 = \emptyset$ and $E_1 \cup E_2 \subseteq A \cup B$. 
    Thus $\Jordan{E_1} + \Jordan{E_2} = \Jordan{E_1 \cup E_2} \leq \JordanInner{A \cup B}$. 
    Taking suprema over $E_1$ and $E_2$ gives 
    $\JordanInner{A} + \JordanInner{B} \leq \JordanInner{A \cup B}$.
\end{proof}

\medskip
\begin{proposition}[Properties of Jordan Contents]
    For bounded sets $A, B \subseteq \bb{R}^n$:
    \begin{enumerate}
        \item $\JordanInner{A} \leq \JordanOuter{A}$
    \end{enumerate}
\end{proposition}

\begin{proof}
    For any $E_1 \subseteq A \subseteq E_2$ with $E_1, E_2 \in \mathcal{E}$, we have 
    $\Jordan{E_1} \leq \Jordan{E_2}$ (monotonicity of Jordan content). 
    Taking supremum over $E_1$ and infimum over $E_2$ gives 
    $\JordanInner{A} \leq \JordanOuter{A}$.
\end{proof}

\bigskip
\begin{definition}[Jordan Measurable]
    A bounded set $A$ is \textbf{Jordan measurable} if $\JordanInner{A} = \JordanOuter{A}$. 
    The common value is the \textbf{Jordan content} $\Jordan{A}$.
\end{definition}

\medskip
\begin{theorem}[Characterization of Jordan Measurability]
    A bounded set $A$ is Jordan measurable if and only if for every $\epsilon > 0$, 
    there exist elementary sets $E_1 \subseteq A \subseteq E_2$ with 
    $\Jordan{E_2 \setminus E_1} < \epsilon$.
\end{theorem}

\begin{proof}
    ($\Rightarrow$) Given $\epsilon > 0$, by definition of sup and inf, there exist 
    $E_1 \subseteq A$ with $\Jordan{E_1} > \JordanInner{A} - \epsilon/2$ and 
    $E_2 \supseteq A$ with $\Jordan{E_2} < \JordanOuter{A} + \epsilon/2$. 
    Since $\JordanInner{A} = \JordanOuter{A}$:
    \[
        \Jordan{E_2 \setminus E_1} = \Jordan{E_2} - \Jordan{E_1} < \epsilon
    \]

    ($\Leftarrow$) We have $\JordanInner{A} \geq \Jordan{E_1}$ and 
    $\JordanOuter{A} \leq \Jordan{E_2}$. Thus:
    \[
        \JordanOuter{A} - \JordanInner{A} \leq \Jordan{E_2} - \Jordan{E_1} 
        = \Jordan{E_2 \setminus E_1} < \epsilon
    \]
    Since $\epsilon$ is arbitrary, $\JordanInner{A} = \JordanOuter{A}$.
\end{proof}

\bigskip
\begin{theorem}[Boundary Characterization]
    A bounded set $A$ is Jordan measurable if and only if $\JordanOuter{\Boundary{A}} = 0$.
\end{theorem}

\begin{proof}
    Note that $\Boundary{A} = \Closure{A} \setminus \Interior{A}$, where $\Closure{A}$ 
    is the closure and $\Interior{A}$ is the interior.

    \medskip
    ($\Rightarrow$) Assume $A$ is Jordan measurable. For any $\epsilon > 0$, there 
    exist elementary sets $E_1 \subseteq A \subseteq E_2$ with 
    $\Jordan{E_2 \setminus E_1} < \epsilon$.

    Since $E_1 \subseteq A$, we have $\Interior{E_1} \subseteq \Interior{A}$. 
    Since $A \subseteq E_2$, we have $\Closure{A} \subseteq \Closure{E_2}$.

    Thus $\Boundary{A} \subseteq \Closure{E_2} \setminus \Interior{E_1} \subseteq E_2 \setminus E_1$ 
    (for appropriate elementary sets).

    So $\JordanOuter{\Boundary{A}} \leq \Jordan{E_2 \setminus E_1} < \epsilon$. 
    Since $\epsilon$ is arbitrary, $\JordanOuter{\Boundary{A}} = 0$.

    \medskip
    ($\Leftarrow$) Assume $\JordanOuter{\Boundary{A}} = 0$. For any $\epsilon > 0$, 
    cover $\Boundary{A}$ by an elementary set $E$ with $\Jordan{E} < \epsilon$.

    Take $E_1 = \Closure{\Interior{A} \setminus E}$ (elementary approximation of interior) 
    and $E_2 = \Closure{A} \cup E$.

    Then $E_1 \subseteq A \subseteq E_2$ and $\Jordan{E_2 \setminus E_1} \leq \Jordan{E} < \epsilon$.
\end{proof}


\bigskip
\begin{proposition}[Jordan Measurable Sets Form an Algebra]
The collection $\mathcal{J}$ of Jordan measurable bounded sets forms an algebra on $\bb{R}^n$ (within any fixed bounding rectangle).
\end{proposition}

\begin{proof}
\textbf{Contains $\emptyset$:} $\JordanInner{\emptyset} = 0 = \JordanOuter{\emptyset}$, so $\emptyset \in \mathcal{J}$.

\medskip
\textbf{Closure under complements:} Let $A \in \mathcal{J}$ be bounded within a rectangle $R$. Then $A^c \cap R$ is bounded. We claim $A^c \cap R \in \mathcal{J}$.

For any $\epsilon > 0$, there exist elementary sets $E_1 \subseteq A \subseteq E_2$ with $\Jordan{E_2 \setminus E_1} < \epsilon$. Then $E_2^c \subseteq A^c \subseteq E_1^c$, so:
\[ \JordanOuter{A^c \cap R} - \JordanInner{A^c \cap R} \leq \Jordan{(E_1^c \cap R) \setminus (E_2^c \cap R)} = \Jordan{E_2 \setminus E_1} < \epsilon \]
Since $\epsilon$ is arbitrary, $A^c \cap R \in \mathcal{J}$.

\medskip
\textbf{Closure under finite unions:} Let $A, B \in \mathcal{J}$. We use the boundary characterization. Since $\Boundary{A \cup B} \subseteq \Boundary{A} \cup \Boundary{B}$:
\[ \JordanOuter{\Boundary{A \cup B}} \leq \JordanOuter{\Boundary{A}} + \JordanOuter{\Boundary{B}} = 0 + 0 = 0 \]
Thus $A \cup B \in \mathcal{J}$.
\end{proof}

\bigskip
\begin{proposition}[Null Sets are Jordan Measurable]
If $\JordanOuter{A} = 0$, then $A$ is Jordan measurable with $\Jordan{A} = 0$.
\end{proposition}

\begin{proof}
We have $0 \leq \JordanInner{A} \leq \JordanOuter{A} = 0$. Thus $\JordanInner{A} = \JordanOuter{A} = 0$, so $A$ is Jordan measurable with $\Jordan{A} = 0$.
\end{proof}

\bigskip
\begin{corollary}[Subsets of Null Sets]
If $\JordanOuter{B} = 0$ and $A \subseteq B$, then $A$ is Jordan measurable with $\Jordan{A} = 0$.
\end{corollary}

\begin{proof}
By monotonicity of outer content, $\JordanOuter{A} \leq \JordanOuter{B} = 0$. Thus $\JordanOuter{A} = 0$, and by the previous proposition, $A$ is Jordan measurable with $\Jordan{A} = 0$.
\end{proof}

\bigskip
\begin{compendium}

\textbf{Contentable Spaces:} A pair $(X, \mathcal{A})$ where $\mathcal{A}$ is an algebra on $X$.

\textbf{Content:} A function $\mu: \mathcal{A} \to [0, \infty]$ with $\mu(\emptyset) = 0$ and finite additivity.

\textbf{Content Properties:}
\begin{itemize}
    \item Monotonicity: $A \subseteq B \implies \mu(A) \leq \mu(B)$
    \item Subtractivity: $A \subseteq B$, $\mu(A) < \infty \implies \mu(B \setminus A) = \mu(B) - \mu(A)$
    \item Finite subadditivity: $\mu(A \cup B) \leq \mu(A) + \mu(B)$
\end{itemize}

\textbf{Outer Content Axioms (O1--O3):}
\begin{enumerate}
    \item[(O1)] $\mu^*(\emptyset) = 0$
    \item[(O2)] Monotonicity
    \item[(O3)] Finite subadditivity
\end{enumerate}

\textbf{Inner Content Axioms (I1--I3):}
\begin{enumerate}
    \item[(I1)] $\mu_*(\emptyset) = 0$
    \item[(I2)] Monotonicity
    \item[(I3)] Superadditivity on disjoint sets
\end{enumerate}

\textbf{Jordan Content:}
\begin{itemize}
    \item Outer: $\JordanOuter{A} = \inf\{\Jordan{E} \mid A \subseteq E \in \mathcal{E}\}$
    \item Inner: $\JordanInner{A} = \sup\{\Jordan{E} \mid E \subseteq A, E \in \mathcal{E}\}$
    \item Measurable: $\JordanInner{A} = \JordanOuter{A}$
\end{itemize}

\textbf{Key Theorems:}
\begin{itemize}
    \item $A$ is Jordan measurable $\iff$ $\JordanOuter{\Boundary{A}} = 0$
    \item Jordan measurable sets form an algebra
    \item Null sets and their subsets are Jordan measurable
\end{itemize}

\end{compendium}

